\section{A common periodic set and consequences}
\label{sec:global}

\subsection{Only contracting copies of the periodic sets are needed}

We prove \cref{thm:main}. For $j\ge1$ put
\begin{equation}\label{eq:ratio-exhaustion}
 I_j=\left[\frac1{j+2},1-\frac1{j+2}\right],
 \qquad \bigcup_{j\ge1}I_j=(0,1).
\end{equation}
For $j\ge1$ and $k\ge0$, use \cref{prop:normalized} on $I_j$ with
\begin{equation}\label{eq:global-budget}
 p_{j,k}=\frac\eps{24}\,2^{-j}2^{-(k+1)}
\end{equation}
to obtain $H_{j,k}$. Define
\begin{equation}\label{eq:global-C}
 C=\bigcup_{j\ge1,\ k\ge0}
 \left(2^{-k}H_{j,k}\ \cup\ (-2^{-k}H_{j,k})\right),
 \qquad F_\eps=\R\setminus C.
\end{equation}
Every component in this union is open with period $2^{-k}$ and hence
also period one. The union is open, symmetric, and $1$-periodic.
The density of a contracted component in a unit interval is
\begin{equation}\label{eq:contraction-density}
 \rho(2^{-k}H_{j,k})=\rho(H_{j,k}).
\end{equation}
It is \emph{not} $2^{-k}\rho(H_{j,k})$: there are $2^k$ complete
periods in a unit interval. The summable budget explicitly pays for
this fact. Reflections have the same density, and therefore
\begin{equation}\label{eq:global-density}
 \rho(C)\le12\sum_{j\ge1,k\ge0}p_{j,k}=\frac\eps2.
\end{equation}
Periodicity gives \eqref{eq:main-density} for every $u$.

Fix $x$, $s\ne0$, and $q\in(0,1)$, and choose $j$ with $q\in I_j$.
Choose an integer $h\ge0$ so that $0<|s|q^h\le1$.
There are $k\ge0$ and $t\in[1,2)$ with
\begin{equation}\label{eq:tail-normalization}
 |s|q^h=2^{-k}t.
\end{equation}
If $s>0$, apply the normalized hitting property of $H_{j,k}$ at center
$2^kx$. Some $n\ge1$ satisfies
\[
 2^kx+tq^n\in H_{j,k},
 \quad\text{so}\quad x+sq^{h+n}\in2^{-k}H_{j,k}\subset C.
\]
If $s<0$, apply it at $-2^kx$ and reflect; the hit lies in
$-2^{-k}H_{j,k}$. This proves that every orbit meets $C$.
Apply the same assertion after any initial tail, replacing $s$ by
$sq^J$, to obtain arbitrarily late hits. The proof of
\cref{thm:main} is complete.

This contraction-only assembly improves the global organization of
the fixed-ratio construction. It uses the elementary identity that a
tail of $G_q$ is a scaled copy of $G_q$. It does not require a relation
between the ratio $q$ and the dyadic grid base.

\subsection{Large compact subsets of arbitrary measurable sets}

\begin{corollary}[Avoidance inside any measurable set]\label{cor:inside}
If $A\subset\R$ is measurable with $0<m(A)<\infty$, then for every
$\delta>0$ there is a compact $K\subset A$ such that
$m(K)>m(A)-\delta$ and $K$ contains no affine copy of any $G_q$,
$0<q<1$.
\end{corollary}

\begin{proof}
Choose compact $A_0\subset A$ with $m(A_0)>m(A)-\delta/2$ and a
bounded interval $I$ containing $A_0$. An affine image of the compact
set in \cref{thm:main}, with sufficiently small accuracy, gives a
compact $K_0\subset I$ with $m(I\setminus K_0)<\delta/2$.
Set $K=A_0\cap K_0$. Avoidance is inherited by subsets and preserved
by invertible affine maps; the measure bound follows.
\end{proof}

\begin{corollary}[A fat Cantor avoiding set]\label{cor:fat-cantor}
The compact set in \cref{thm:main} can be chosen perfect and nowhere
dense without losing measure. In particular there is a fat Cantor
set answering \cite[Question~1]{BGKMW} negatively.
\end{corollary}

\begin{proof}
The set $F_\eps$ has empty interior, since every nonempty interval
contains an affine contracting geometric progression. Its compact
section $K$ is therefore nowhere dense. Let $K'$ be the support of
the finite measure $\ind_K\,dm$. It is a closed subset of $K$;
its complement in $K$ is covered by countably many rational intervals
of that measure zero, so $m(K')=m(K)$. The support has no isolated
points: a measure supported on one isolated point of a neighborhood
would have an atom there, whereas $\ind_K\,dm$ is atomless.
Thus $K'$ is perfect and has all the required properties.
\end{proof}

The effective-closedness claim later concerns the original construction.
We do not infer an effective procedure for taking the measure support
in \cref{cor:fat-cantor}.

\begin{corollary}[Finite-dimensional line patterns]\label{cor:dimensions}
For each $d\ge1$ and $\delta>0$, there is a compact
$E\subset[0,1]^d$ of $d$-dimensional measure greater than $1-\delta$
containing no set $\{x+q^nv:n\ge1\}$ with $x\in\R^d$,
$v\in\R^d\setminus\{0\}$, and $0<q<1$.
\end{corollary}

\begin{proof}
Take $E=K^d$ with $m(K)$ sufficiently close to one. A nonzero $v$
has a nonzero coordinate. Projection to that coordinate would give
an affine copy of $G_q$ in $K$ if the whole pattern lay in $E$.
\end{proof}

The periodic set in \cref{thm:main} is asserted to avoid contracting
ratios. A nonempty periodic set cannot exclude all expanding geometric
orbits: it contains $x+2^n$ whenever it contains $x$. The compact
section automatically excludes expanding orbits because they are
unbounded. The range $0<q<1$ keeps the main statement precise.
